\uuid{0w28}
\titre{DL usuels/Calcul de limites/Étude locale}
\auteur{hanani ; mzari}
\organisation{exo7}
\licence{CC BY-NC-SA 4.0}
\datecreate{2020-01-20}
\niveau{L1}
\module{Analyse}
\chapitre{Développement limité}
\sousChapitre{Applications}
\difficulte{3}
\theme{Applications des DL, DL usuels, Calcul de limites, Étude locale}
\exo7qcmid{lille-2.developpements-limites.041}

\contenu{
\texte{
Soit $a\in \Rr$ et $f$ une fonction. On dit que $f(x)$ est un $o\left((x-a)^n\right)$ au voisinage de $a$ et on écrit $f(x)=o((x-a)^n)$ si $\displaystyle \lim _{x\to a}\frac{f(x)}{(x-a)^n}=0$. On note par $DL_n(a)f(x)$ le développement limité de $f$ à l'ordre $n$ en $a$. On note aussi $DL_n(a^+)f(x)$ (resp. $DL_n(a^-)f(x)$) le développement limité de $f$ à l'ordre $n$ à droite en $a$ (resp. à gauche en $a$).
}

\begin{question}
\qclassification{125.03}
\qtags{motcle=DL usuels/Calcul de limites/Étude locale}
Soit $\displaystyle f(x)=\sqrt{1+2x}-\sqrt[3]{1+3x}$ et $g(x)=1-\cos x$. On note $T_0$ la tangente au graphe de $f$ au point $0$. Parmi les affirmations suivantes, cocher celles qui sont vraies :
\begin{answers}
\bad{$T_0$ est horizontale et le graphe de $f$ est en dessous de $T_0$ au voisinage de $0$.}
\good{$\displaystyle DL_2(0)g(x)=\frac{x^2}{2}+o(x^2)$.}
\bad{$\displaystyle DL_0(0)\frac{f(x)}{g(x)}=0+o(1)$.}
\good{$\displaystyle \lim _{x\to 0}\frac{\sqrt{1+2x}-\sqrt[3]{1+3x}}{1-\cos x}=1$.}
\end{answers}
\vskip2mm
\begin{explanations}
Les développements à l'ordre $2$ en $0$ de $(1+t)^{\alpha}$ et de $\cos (x)$ donnent :
$$\displaystyle DL_2(0)f(x)=\frac{x^2}{2}+o(x^2)\quad \mbox{et}\quad DL_2(0)g(x)=\frac{x^2}{2}+o(x^2).$$
D'abord, $\displaystyle \lim _{x\to 0}\frac{\sqrt{1+2x}-\sqrt[3]{1+3x}}{1-\cos x}=1$. Ensuite, $T_0$ est la droite d'équation $y=0$. Or, au voisinage de $0$, on a : $\displaystyle f(x)-y\simeq \frac{x^2}{2}\geq 0$. Donc le graphe de $f$ est situé au dessus de $T_0$.
\end{explanations}
\end{question}
}
